% 2013 年北京卷文科第 10 题：四棱锥的三视图。
% 来源 PDF：普通高考/2013/2013北京文.pdf，第 1 页（印面 789）中栏。
% 正视图底边 2+1=3、高 1；侧视图底边 1+2=3、高 1；俯视图为边长 3 的正方形，
% 内部四条线汇聚于顶点的投影（距左边 2、距上边 1）。
\begin{tikzpicture}[font=\footnotesize, line width=0.45pt, >=Stealth]
  % 正（主）视图
  \draw (0,0) -- (3,0) -- (2,1) -- cycle;
  \draw[|<->|] (0,-0.30) -- (1.95,-0.30) node[midway,below] {$2$};
  \draw[|<->|] (2.05,-0.30) -- (3,-0.30) node[midway,below] {$1$};
  \foreach \x in {0,2,3} \draw (\x,-0.18) -- (\x,-0.42);
  \node at (1.5,-1.0) {正（主）视图};
  % 高 1（正视图与侧视图之间）
  \draw[gray!70] (2,1) -- (3.55,1);
  \draw[gray!70] (3,0) -- (3.55,0);
  \draw[|<->|] (3.42,0) -- (3.42,1) node[midway,left] {$1$};
  % 侧（左）视图
  \draw (4.5,0) -- (7.5,0) -- (5.5,1) -- cycle;
  \draw[|<->|] (4.5,-0.30) -- (5.45,-0.30) node[midway,below] {$1$};
  \draw[|<->|] (5.55,-0.30) -- (7.5,-0.30) node[midway,below] {$2$};
  \foreach \x in {4.5,5.5,7.5} \draw (\x,-0.18) -- (\x,-0.42);
  \node at (6.0,-1.0) {侧（左）视图};
  % 俯视图
  \draw (1.2,-1.7) rectangle (4.2,-4.7);
  \coordinate (Q) at (3.2,-2.7);
  \draw (1.2,-1.7) -- (Q);
  \draw (4.2,-1.7) -- (Q);
  \draw (1.2,-4.7) -- (Q);
  \draw (4.2,-4.7) -- (Q);
  \node at (2.7,-5.15) {俯视图};
\end{tikzpicture}
